Sections~\ref{sec:cover} and~\ref{sec:multi} show that a multi-reference plan
depends on how the latent spread is divided between a ship-wide term $\tau$ and
a ship--product term $\omega$. Whether a fleet's records could estimate that
division, and what a plan built from them delivers, can be answered by
simulation before any record is collected. Every number in this section is
simulated under assumed parameters.

\subsection{Design of the Study}
For ship $s$, reference $p$ and mission $t$ the generating model is
\begin{equation}
y_{spt}\sim\mathrm{Poisson}\!\left(m\,a_p\exp\{u_s+w_{sp}+e_{spt}
-\tfrac12(\tau^2+\omega^2+\varepsilon^2)\}\right),
\label{eq:sim}
\end{equation}
with $u_s\sim\mathcal N(0,\tau^2)$, $w_{sp}\sim\mathcal N(0,\omega^2)$ and a
mission term $e_{spt}\sim\mathcal N(0,\varepsilon^2)$, so that each cell expects
$m\,a_p$ failures. The reference weights $a_p$ are equal or proportional to
$1:2:4:8$. Five allocations are studied: ship-wide only $(\tau,\omega)=(0.6,0)$,
interaction only $(0,0.6)$, equal splits $(0.42,0.42)$ and $(0.3,0.3)$, and none
$(0,0)$. A full factorial in the number of ships $S\in\{5,10,20,40,80\}$ and
missions per ship $M\in\{1,2,4,8\}$, with $P=4$, $m=1$ and $\varepsilon=0.3$, is
completed by one-factor variations of $m\in\{0.1,0.3,3\}$, $P\in\{2,8\}$,
$\varepsilon\in\{0,0.6\}$ and unequal weights: \dsCells{} cells of \dsReps{}
campaigns each, \dsCampaigns{} in all.

Every campaign is fitted three times with the estimator of
Section~\ref{sec:model} (both terms, ship term only, interaction term only,
reference as a fixed effect), which gives the estimates and the likelihood-ratio
statistic for each term. Because a variance at zero lies on the boundary of the
parameter space, the statistic is referred to an equal mixture of
$\chi^2_0$ and $\chi^2_1$~\cite{self1987,stram1994}. In 15 cells the first 300
campaigns also receive profile-likelihood intervals for $\tau$ and $\omega$.
A 99\,\% plan for one mission of an unseen ship is then built from each fit in
four ways: \emph{plug-in}, from draws of the latent terms at the estimated
parameters; \emph{predictive}, adding draws of the coefficients from their
approximate posterior; \emph{conservative}, the larger of the plug-in plans with
the estimated total spread put wholly on $\tau$ and wholly on $\omega$; and
\emph{Poisson}, each reference at its own Poisson quantile. The cover each plan
delivers is computed exactly under the true model~\eqref{eq:sim} by
Proposition~\ref{prop:multi}, with the mission term folded into the reference
term. Monte Carlo standard errors follow~\cite{morris2019}: \dsMcseRej{}
percentage points for a rejection rate near 5\,\%.

No campaign failed to fit. The conditioning is specific: the latent effects are
Gaussian, $\varepsilon$ is fixed within a cell and there are no common-cause
events, reference changes or reporting errors. Section~\ref{sec:fleet} relaxes
those assumptions.

\subsection{Testing and Estimating the Allocation}
The boundary test holds its size. Across the cells where a component is zero
the rejection rate at the 5\,\% level has median \dsTypeOneMed\,\% and maximum
\dsTypeOneMax\,\%, conservative as the mixture theory predicts. A spurious
component is therefore no more frequent than the test allows, provided it is
tested and not judged from whether a point estimate leaves the boundary, which
happens about half the time under the null whatever the campaign size.

Table~\ref{tab:power} gives the power under the equal split $(0.42,0.42)$.
With one mission per ship the interaction term is never detected, at any
number of ships: it varies at the same level as the mission term and cannot be
told apart from it, and when the truth is interaction only the larger component
is identified correctly in only \dsOrderOneLo--\dsOrderOneHi\,\% of campaigns.
Repeated missions are what identifies it. With four missions, the interaction
is detected in \dsPowOmTenFour\,\% of ten-ship campaigns and
\dsPowOmTwentyFour\,\% of twenty-ship ones, and with eight missions and forty
ships in \dsPowOmFortyEight\,\%. With at least four missions and ten ships the
larger component is identified in at least \dsOrderFour\,\% of campaigns. The
profile intervals cover the true values in \dsCiLo--\dsCiHi\,\% of the
\dsCiReps{} campaigns per cell against a nominal 95\,\%.

\begin{table}[!t]
\caption{Power of the boundary likelihood-ratio test at 5\,\% when both
components are present, $(\tau,\omega)=(0.42,0.42)$, in percent
(\textsc{simulation}, \dsReps{} campaigns per cell, $P=4$, $m=1$)}
\label{tab:power}
\centering
\footnotesize
\setlength{\tabcolsep}{4.5pt}
\begin{tabular}{@{}ccccccccc@{}}
\toprule
& \multicolumn{4}{c}{ship-wide $\tau$, missions $M$} &
\multicolumn{4}{c}{interaction $\omega$, missions $M$} \\
\cmidrule(lr){2-5}\cmidrule(lr){6-9}
$S$ & 1 & 2 & 4 & 8 & 1 & 2 & 4 & 8 \\
\midrule
\tabinput{tables/tab_power.tex}
\bottomrule
\end{tabular}
\end{table}

\subsection{Cover of Plans for an Unseen Ship}
Table~\ref{tab:campaign} reports the joint cover delivered by each plan,
averaged over the five allocations. The predictive plan reaches its target from
ten ships on (\dsPredCoverTen\,\%), and in no cell with at least ten ships and
two missions does its mean fall below \dsPredMin\,\%. The plug-in plan, which
ignores the uncertainty of the coefficients, falls short when the campaign is
small (\dsPlugCoverFive\,\% with five ships). The conservative plan behaves
like the plug-in one at slightly higher cost. The Poisson plan delivers
\dsPoisLo--\dsPoisHi\,\% across allocations, \dsPoisMThree\,\% when each cell
expects three failures and \dsPoisPEight\,\% with eight references, since its
shortfall grows with the count and with the number of references that must
all be covered.

The price of a small campaign shows in the stock rather than in the cover.
With five ships the predictive plan has a median of \dsUnitsFiveMed{} units,
close to the \dsUnitsOracle{} of the plan that knows the parameters, but a mean
of \dsUnitsFiveMean: in a minority of campaigns the coefficients are so poorly
determined that the plan becomes very large. That is the correct response of a
predictive distribution to little evidence, and it is what a fleet with five
ships should expect to pay.

\begin{table}[!t]
\caption{Mean joint cover of a 99\,\% plan for a four-reference zone of an
unseen ship, in percent, averaged over five allocations, with the mean tenth
percentile of the predictive plan and median units of the predictive and
oracle plans (\textsc{simulation}; $M=4$)}
\label{tab:campaign}
\centering
\footnotesize
\setlength{\tabcolsep}{3.2pt}
\begin{tabular}{@{}cccccccc@{}}
\toprule
$S$ & plug-in & predictive & conservative & Poisson & pred. p10 &
\multicolumn{2}{c}{units (pred./oracle)} \\
\midrule
\tabinput{tables/tab_campaign_cover.tex}
\bottomrule
\end{tabular}
\end{table}

\subsection{Implications for a First Campaign}
Within the configurations examined, a campaign should follow several references
on the same ships over repeated missions; breadth in ships does not substitute
for repetition, and a cross-sectional snapshot cannot separate the two terms at
all. Ten ships over four missions is the smallest configuration in this grid
that orders the components reliably, and twenty to forty ships detect the
interaction with high probability. Plans should be computed from the full
predictive distribution, which keeps its cover in small campaigns at the cost
of stock, and the components should be reported with their intervals and
tests.
